% MAIN EDITING SPACE FOR FUTURE ECOLOGICAL RESULTS.
% Replace pending text with measured findings; never replace a dash by a guess.
% Keep config, revision, seed list, selection provenance and raw result paths.
% Add figures under figures/ as PDF/PNG/JPG and use relative paths.
\section{Ecological Observations and Experimental Results}\label{sec:observations}
\pending{Reserved for future experiments. No ecological learning outcomes have
yet been entered into this manuscript. The engineering measurements in
the engineering Results are historical system measurements and are not
substitutes for the studies below.}

\subsection{Experiment registry}
For each completed study, record the implementation revision, environment JSON,
map, algorithm per species, observation version/$k$, architecture, population
settings, seed policy/list, step/update/trial budgets, reproductive/mutation
settings, champion criterion and initialization source. Preserve the effective
configuration and distinguish fresh, warm-start and exact-resume conditions.

\begin{table}[ht]
\centering\small
\caption{Fillable registry. Entries are proposed study families, not completed
experiments. Replace the placeholders with actual identifiers and settings.}
\label{tab:registry}
\begin{tabularx}{\linewidth}{@{}lYYY@{}}\toprule
Study & Main comparison & Run IDs / config / revision & Seeds and budget\\\midrule
S1 & Trained / fresh / random & \pending{To add} & \pending{To add}\\
S2 & Sensing/capacity ablation & \pending{To add} & \pending{To add}\\
S3 & Inheritance/mutation ablation & \pending{To add} & \pending{To add}\\
S4 & Geography transfer & \pending{To add} & \pending{To add}\\
S5 & Champion-selection criterion & \pending{To add} & \pending{To add}\\\bottomrule
\end{tabularx}
\end{table}

\subsection{Survival, feeding and hunting}
\pending{Insert population trajectories, extinction outcomes, completed-lifespan
and censored-survivor summaries, feeding/hunting rates and energy coverage.
Describe whether differences persist on held-out seeds with frozen opponents.
Report denominators and uncertainty; separate larger populations from higher
per-animal success.}

\begin{table}[ht]
\centering\small
\caption{Outcome template. Dashes mean no measurement supplied; they are not
zeros. Split by species, terrain and comparison arm as needed.}
\label{tab:future-outcomes}
\begin{tabularx}{\linewidth}{@{}Yccc@{}}\toprule
Endpoint & Trained & Fresh & Random\\\midrule
Extinction fraction across seeded runs & --- & --- & ---\\
Completed lifespan / survivor censoring & --- & --- & ---\\
Kills per 1,000 predator animal-steps & --- & --- & ---\\
Successful feeds / feed attempts & --- & --- & ---\\
Food-energy intake / measured energy demand & --- & --- & ---\\
Births and viable descendants & --- & --- & ---\\\bottomrule
\end{tabularx}
\end{table}

\subsection{Reproduction, mutation and family selection}
\pending{Insert paired reproduction attempt/success rates, parent energy costs,
newborn grants, generations and family statistics. Compare crossover and mutation
conditions under matched budgets. Distinguish inherited performance from gains
learned by descendants. Report whether lineage selection transfers to new
worlds rather than only increasing its training-time ranking score.}

\subsection{Geography and observation ablations}
\pending{Insert swim attempts/success, drowning, current-driven movement,
climbing demand and map-specific survival/reproduction. Compare observation
widths or sight/capacity changes with compatible fresh-trained models. Identify
whether improvements concern perception, policy optimization or physical tuning.}

\subsection{Learning signals, uncertainty and failure cases}
\pending{Insert update/entropy trends and GRPO return spread alongside ecological
outcomes. Report equal-return trial fractions, missing diagnostic coverage,
failed mating choices and representative unsuccessful seeds. Include confidence
intervals and the analysis procedure. A visually appealing replay alone is not
an experimental result.}

% Example figure once real data exist (remove comments and supply the file):
% \begin{figure}[t]
% \centering
% \includegraphics[width=.9\linewidth]{figures/held-out-populations.pdf}
% \caption{Population trajectories across N held-out seeds. Lines show ...;
% bands show ... . Models were frozen and selected using ... .}
% \label{fig:held-out-populations}
% \end{figure}
